% Generated draft: requires mathematical and source review.
\documentclass[11pt]{article}
\usepackage[margin=25mm]{geometry}
\usepackage{fontspec}
\setmainfont{Latin Modern Roman}
\usepackage{amsmath,amssymb,mathtools,unicode-math}
\setmathfont{Latin Modern Math}
\usepackage{enumitem,needspace,xurl,hyperref}
\setlength{\parindent}{0pt}
\setlength{\parskip}{5pt}
\setlength{\emergencystretch}{3em}
\newcommand{\N}{\mathbb{N}}
\newcommand{\Z}{\mathbb{Z}}
\newcommand{\Q}{\mathbb{Q}}
\newcommand{\R}{\mathbb{R}}
\newcommand{\C}{\mathbb{C}}
\newcommand{\F}{\mathbb{F}}
\newcommand{\A}{\mathbb{A}}
\newcommand{\PP}{\mathbb{P}}
\newcommand{\E}{\mathbb{E}}
\newcommand{\eps}{\varepsilon}
\newcommand{\dd}{\,\mathrm{d}}
\begin{document}
\section*{UnsolvedMath ID: 4600038}
\textbf{Source problem code:} catalogue entry
\subsection{Definitions and mathematical statement}
A two-sided subshift is a closed shift-invariant subset $X\subseteq\mathcal A^{\mathbb Z}$ for a finite alphabet $\mathcal A$, with shift map $T$. It is a shift of finite type (SFT) if finitely many forbidden blocks define it, and is irreducible if any two allowed words can be joined by an allowed connecting word. Two such systems $(X,T)$ and $(Y,S)$ are topologically orbit equivalent in the orbit-homeomorphism sense if a homeomorphism $h:X\to Y$ satisfies
\[
h(\{T^n x:n\in\mathbb Z\})=\{S^n h(x):n\in\mathbb Z\}\qquad(x\in X).
\]
Whether additional continuous orbit-cocycle or orientation requirements are intended is not stated.

For context, the usual integer cohomology group and its standard positive cone are
\[
H(X,T)=C(X,\mathbb Z)/\{g-g\circ T:g\in C(X,\mathbb Z)\},\qquad
H(X,T)^+=\{[f]:f\geq0\},
\]
with distinguished unit $u=[1]$. Here $C(X,\mathbb Z)$ consists of continuous integer-valued functions. An isomorphism of unital ordered groups preserves both the positive cone and $u$.

The source asks for the following related classification problems:
\begin{enumerate}
\item Classify irreducible SFTs up to topological orbit equivalence.
\item Characterize which unital ordered cohomology groups occur for irreducible SFTs, and classify the resulting ordered groups with their units.
\item Determine when the winding order and the standard order on the relevant cohomology agree.
\item If a subshift is topologically orbit equivalent to an irreducible SFT, must it itself be an SFT?
\end{enumerate}
The record does not define winding order or identify the exact cohomology/order conventions, so the usual cohomology construction above explains the setting but cannot supply the missing winding cone. Existing ordered-dimension-group and flow-equivalence results mentioned in the summary are background, not an assertion that flow equivalence and orbit equivalence coincide.
\subsection{Short English statement}
Classify irreducible shifts of finite type by topological orbit equivalence and understand their ordered cohomology groups. Compare winding and standard orders, and ask whether orbit equivalence to an irreducible SFT forces a subshift to be an SFT.
\subsection{Sources}
\begin{itemize}
\item Boyle, \emph{Open Problems in Symbolic Dynamics} (2008), item 23.4, \url{https://www.math.umd.edu/~mboyle/papers/openfinalsub3nov2007.pdf}; accessed 2026-07-29, with PDF-text extraction recorded in the supplied provenance.
\item The supplied research summary attributes developed ordered-dimension-group/flow-equivalence theory to Boyle--Handelman, without a separate publication citation.
\item Local registry \texttt{data/unsolvedmath/problems.json}, permanent UnsolvedMath ID $4600038$, source code catalogue entry, \emph{Orbit equivalence and ordered cohomology}. Dataset attribution: Ulam AI and the UnsolvedMath Contributors.
\end{itemize}
\end{document}
